\pdfinfoomitdate=1
\pdftrailerid{}
\documentclass[10pt]{article}
\usepackage[a4paper,margin=0.70in]{geometry}
\usepackage[T1]{fontenc}
\usepackage{lmodern}
\usepackage{amsmath,amssymb,booktabs,microtype}
\title{Exact Two-Clock Periodic Geometry of Cyclic Rule 90}
\author{Route-A structural certificate C145}
\date{}
\begin{document}
\maketitle

\begin{abstract}
We count every temporally periodic configuration of Rule 90 on every finite
cyclic binary lattice.  For spatial circumference $L$ and time $n$, the fixed
count is a power of two whose exponent is an explicit polynomial-gcd degree.
An ideal-theoretic proof covers even $L$, where the spatial polynomial can be
non-squarefree.  M\"obius inversion then separates exact-period points from
primitive temporal cycles, and fixed configurations are exactly labeled
$L\times n$ spatiotemporal tori.  A complete $24\times24$ ledger gives
domain-qualified aspect-ratio witnesses and a same-fixed-count control with
different primitive content.  The result is exact finite-volume progress, not
a single frozen target determinant.
\end{abstract}

\section{Frozen Rule-90 family}
On the cyclic lattice $\mathbb Z/L\mathbb Z$, Rule 90 is
\[
(F_Lu)_i=u_{i-1}+u_{i+1}\pmod2.
\]
Both summands are retained when the two neighbor indices coincide; over
$\mathbb F_2$ they then cancel.  Thus the definition also covers $L=1,2$.
Encode $u$ in
$R_L=\mathbb F_2[x,x^{-1}]/(x^L-1)$; then $F_L$ is multiplication by
$a=x+x^{-1}$.  Freeze the ordered pair $(L,n)$ as the spatial and temporal
clocks.  Write $\operatorname{Fix}(L,n)=\#\ker(F_L^n-I)$,
$P_L(n)$ for points of least positive temporal period $n$, and $C_L(n)$ for
their geometric temporal cycles.  These three quantities are not
interchangeable.

\section{Kernel--gcd theorem}
\textbf{Multiplication lemma.}  For monic $f\in k[x]$ and any $h$, multiplication
$m_h$ on $k[x]/(f)$ satisfies
\begin{equation}
\dim\ker m_h=\deg\gcd(f,h).                                      \tag{1}
\end{equation}
Indeed, put $g=\gcd(f,h)$, $f=gf_1$, and $h=gh_1$.  Since
$\gcd(f_1,h_1)=1$,
\[
f\mid hq\quad\Longleftrightarrow\quad f_1\mid q.
\]
The kernel is therefore represented uniquely by $f_1r$ with $\deg r<\deg g$,
which proves (1): these representatives span, and
$f\mid f_1(r-r')$ would force $g\mid(r-r')$, hence $r=r'$ by the degree
bound.  This also covers $\deg g=0$.  The argument uses divisibility rather
than distinct roots.

Because $x$ is invertible modulo $x^L+1=x^L-1$ in characteristic two,
clearing the Laurent denominator does not change the kernel:
\begin{equation}
x^n(a^n-1)=(x^2+1)^n+x^n=:b_n.                                  \tag{2}
\end{equation}
Applying (1) gives the all-size formula
\begin{equation}
\boxed{\ \operatorname{Fix}(L,n)
=2^{\deg\gcd(x^L+1,(x^2+1)^n+x^n)}\ },\qquad L,n\geq1.           \tag{3}
\end{equation}
No squarefree hypothesis occurs: for example
$x^6+1=(x^3+1)^2$ over $\mathbb F_2$, and (3) still applies.

\section{Exact temporal cycles and tori}
Even when $F_L$ is not invertible on the whole state space, a point satisfying
$F_L^nu=u$ lies on a cycle whose least period divides $n$.  Hence the
fixed-point set at time $n$ is the disjoint union of exact-period sets at
divisors of $n$.  M\"obius inversion and orbit partitioning give
\begin{equation}
P_L(n)=\sum_{d\mid n}\mu(n/d)\operatorname{Fix}(L,d),
\qquad C_L(n)=\frac{P_L(n)}{n}\in\mathbb Z_{\geq0}.               \tag{4}
\end{equation}
The division occurs only after inversion: exact-period-$n$ points form cycles
of size $n$.

A labeled spatiotemporal torus is an array satisfying
\[
u_{i,j+1}=u_{i-1,j}+u_{i+1,j},\qquad i\bmod L,\quad j\bmod n.
\]
Its row at $j=0$ determines all rows, and temporal closure is precisely
$F_L^nu=u$.  Thus the number of labeled $L\times n$ tori is
$\operatorname{Fix}(L,n)$.

\section{Two-clock witnesses}
The table records exact counts for $1\leq L,n\leq24$.  Every minimum below is
within that frozen cutoff, not over unbounded $L,n$; the bounded-search domains
are distinct:
\begin{center}
\begin{tabular}{@{}llll@{}}
\toprule
search subset & area & cells $(L,n;\operatorname{Fix})$ & nonzero exact content\\
\midrule
$1\leq L,n\leq24$ & $3$ & $(1,3;1),(3,1;4)$ & $(3,1):P=4$\\
$2\leq L,n\leq24$ & $6$ & $(2,3;1),(3,2;4)$ & none at displayed times\\
$2\leq L,n\leq24$, some $P>0$ & $12$ & $(2,6;1),(3,4;4),(4,3;1),(6,2;16)$ & $(6,2):P=12,C=6$\\
\bottomrule
\end{tabular}
\end{center}
For each area the search enumerates every eligible divisor pair $(L,n)$ in the
stated box; no best-case pair is omitted.
Thus area $Ln$ loses aspect ratio.  A different control fixes $L=5$:
\[
\operatorname{Fix}(5,3)=\operatorname{Fix}(5,6)=16,
\quad(P_5(3),C_5(3))=(15,5),\quad(P_5(6),C_5(6))=(0,0).           \tag{5}
\]
One fixed count therefore loses the divisor history that determines primitive
content.

\section{Validation and Route-A boundary}
Across 576 cells, the ledger sums to 488,334 fixed points, 283,758
exact-period points, and 24,474 primitive cycles.  The independent
binary-matrix checker passes 6,520 assertions, including direct state
enumeration at small sizes.  SymPy passes 1,177 exact polynomial checks,
production replays byte-identically, and all 43 hostile cases (42 repaired-hash
and one stale-hash) are rejected.

The strict tuple is $(\texttt{A1\_WEAK},\texttt{A2\_FAIL},
\texttt{A3\_FAIL},\texttt{A4\_FAIL})$, overall
\texttt{ROUTE\_A\_EXPLORATORY}.  The intrinsic finite-volume cycles depend on
two essential clocks; we supply no single frozen target divisor, target
functional equation or counting law, arithmetic/local factor, root number,
automorphy, natural operator lift, Hilbert--P\'olya operator, infinite-volume
limit, or Route-B authorization.  The scope literal is
\texttt{NO\_BAD\_EULER\_OR\_ROOT\_NUMBER}.
\end{document}
